\section{Introduction}
\label{sec:intro}

A practitioner building an entity-extraction pipeline needs canonical
identifiers and writes an unremarkable constraint: \idpattern, a
3--31-character lowercase identifier beginning with a letter and
containing letters, digits, or underscores. Asked for the identifier of
an entity whose name is a street address, a GPT-5.6 model in our study
returns \code{888\_7th\_avenue}. The rendering is faithful to the name and
fails the anchor \code{\^{}[a-z]}. A conforming identifier exists
(\code{building\_888\_7th\_avenue}), but the task does not specify which
conforming normalisation to use.

Under the evaluation practice common for structured output, this failure
is indistinguishable from a model that emitted no JSON, a model that
returned the schema definition instead of an instance, and a model that
produced a conformant object inside a markdown fence. All four increment
the same counter, although they have different causes and call for
different remedies.

% NEW: answers "why not constrained decoding / native structured output?"
% before the reader asks it. Only cites what is already in the bib.
Constrained decoding and vendor structured-output modes prevent many
syntactic failures, but they do not remove the need to validate and
repair. Support for JSON Schema features varies between implementations
\citep{geng2025jsonschemabench}, and Pydantic model validators, such as
cross-field agreement and membership checks, have no JSON Schema
counterpart and cannot be enforced at decoding time. A pipeline that
must run across vendors, model generations, or endpoints therefore
keeps the schema in the prompt and validates afterwards, and its
validate-and-retry loop stays in place even where decoding is
constrained. How that loop fails is an engineering question, and a single
validity rate cannot answer it.

We study the prompt-and-validate paradigm: the schema is stated in the
prompt, the completion is validated with Pydantic, and failures are
retried with the validation error appended. We do not evaluate
constrained decoding (\cref{sec:related}). We ask three questions.
\begin{description}[leftmargin=1.4em,itemsep=1pt,topsep=2pt,font=\normalfont\bfseries]
  \item[RQ1] Do models with similar failure rates fail in the same way?
  \item[RQ2] Does a model's failure profile stay stable when the schema's
  constraints or the task domain change?
  \item[RQ3] What do formatting failures conceal, and how does that
  change the reading of repairs and prompt ablations?
\end{description}
We define a taxonomy of conformance failures, assign each failing attempt
a category with a deterministic classifier, and apply it to
\pn{reported.attempts} attempts across \pn{reported.cells} configurations
of \pn{reported.models} models from two vendors, one baseline schema and
two targeted variants, and two task domains.

The answer to all three questions is that the validity rate hides what an
engineer needs. Models with overlapping failure rates fall into four
profiles defined by their dominant category (RQ1). Raising length caps
changes which category dominates for the same model, and the same model
shows different profiles on the two domains (RQ2). Only
\pn{fence.reported.pure_share} of fenced failures validate once the fence
is removed. The rest conceal a length or pattern violation that the
returned parse error never mentions, and ignoring this makes a prompt
ablation appear to move failures from one mode to another (RQ3).

\paragraph{Contributions.}
\begin{enumerate}[leftmargin=1.4em,itemsep=1pt,topsep=2pt]
  \item An operational failure taxonomy with deterministic detectors, an
  explicit precedence rule, and development-set checks against human
  labels (\cref{sec:setup:taxonomy}).
  % DECISION D1: if held-out, second-labeller validation is completed,
  % restore "validated against held-out human labels (Cohen's kappa ...)".
  \item An empirical study of \pn{reported.models} models showing that
  similar failure rates hide different failure modes, and that a profile
  shifts with the constraint surface and the task domain
  (\cref{sec:results:profiles,sec:results:surface}).
  \item A layered analysis showing that markdown wrapping conceals
  underlying schema violations and that ignoring concealment misreads
  both repair loops and prompt ablations (\cref{sec:results:concealment}).
  \item Three practices for engineers who build validation pipelines,
  derived from these results (\cref{sec:implications}).
\end{enumerate}

% EMSE: if review is double-anonymous, use the anonymized mirror for the
% submitted PDF; the arXiv build and the final version use the GitHub URL.
\paragraph{Artifact.} The classifier, harness, locked results database,
raw completions, annotations, and the scripts that regenerate every table
and figure are available at \url{https://github.com/amitvermaknw/schema-enforcement}
(see also \cref{sec:declarations}).
